% Generado por Experimentos/generar_resultados_latex.py desde complementario.
% No editar a mano: volver a ejecutar el script.

\begin{table}[ht]
\centering
\footnotesize
\begin{tabular}{c c c c}
\hline
Ciclo & $\alpha=0$ & $\alpha=0.5$ & $\alpha=1$ \\
\hline
0 ($M_0$) & \multicolumn{3}{c}{$0.0030$ [$0.0012$; $0.0062$]} \\
1 & $0.0053$ [$0.0029$; $0.0070$] & $0.0044$ [$0.0025$; $0.0066$] & $0.0051$ [$0.0025$; $0.0066$] \\
2 & $0.0059$ [$0.0021$; $0.0123$] & $0.0032$ [$0.0004$; $0.0066$] & $0.0140$ [$0.0074$; $0.0206$] \\
3 & $0.0030$ [$0.0012$; $0.0045$] & $0.0043$ [$0.0012$; $0.0066$] & $0.0261$ [$0.0206$; $0.0288$] \\
4 & $0.0047$ [$0.0021$; $0.0070$] & $0.0052$ [$0.0016$; $0.0074$] & $0.0225$ [$0.0198$; $0.0272$] \\
5 & $0.0043$ [$0.0033$; $0.0049$] & $0.0048$ [$0.0025$; $0.0091$] & $0.0350$ [$0.0276$; $0.0387$] \\
\hline
\end{tabular}
\caption{Repetición media de 4-gramas: media y rango [mínimo; máximo] de las tres cadenas por ciclo y proporción sintética $\alpha$. $M_0$ es compartido por las tres proporciones dentro de cada repetición. Salvo la perplejidad, calculada con los 128 contextos de prueba, las métricas corresponden al análisis complementario \textit{post hoc} con los 90 contextos en los que todas las etapas alcanzan $W=30$ palabras.}
\label{tab:medias_repeat_4_mean}
\end{table}
